% 附录A：泛函分析背景
\section{泛函分析背景}
\label{app:functional_analysis}

\subsection{希尔伯特空间}

希尔伯特空间$\mathcal{H}$是完备的内积空间。对于CTUFT，相关的空间是：
\begin{itemize}
    \item Fock空间：$\mathcal{F}_s(L^2(\mathbb{R}^3))$
    \item 单粒子空间：$L^2(\mathbb{R}^3, d^3p/2E_p)$
\end{itemize}

\subsection{算符理论}

\begin{definition}[自伴算符]
算符$A$是自伴的，如果$A = A^\dagger$，即$D(A) = D(A^\dagger)$且对所有$\psi \in D(A)$有$A\psi = A^\dagger\psi$。
\end{definition}

\subsection{谱定理}

\begin{theorem}[谱定理]
对于自伴算符$A$，存在唯一的谱测度$E$使得：
\begin{equation}
    A = \int_{\sigma(A)} \lambda \, dE(\lambda)
\end{equation}
\end{theorem}

\subsection{分布理论}

QFT中的算符值分布满足：
\begin{equation}
    \mathcal{T}(f) = \int d^4x \, \mathcal{T}(x) f(x)
\end{equation}
其中$f \in \mathcal{S}(\mathbb{R}^4)$是测试函数。
